\documentclass[11pt]{article}
\usepackage[margin=1in]{geometry}
\usepackage{amsmath,amssymb,amsthm,mathtools,booktabs,microtype,graphicx,float}
\usepackage[colorlinks=true,linkcolor=blue,urlcolor=blue]{hyperref}
\hypersetup{pdftitle={V1 frequency dispersion and excitation on a square torus}}
\graphicspath{{../simulations/output/v1_extension/}{figures/}}
\newtheorem{proposition}{Proposition}
\newtheorem{corollary}{Corollary}
\newcommand{\ee}{\mathrm e}
\newcommand{\ii}{\mathrm i}
\newcommand{\Rea}{\operatorname{Re}}
\newcommand{\Ima}{\operatorname{Im}}
\newcommand{\diag}{\operatorname{diag}}
\newcommand{\spec}{\operatorname{spec}}
\newcommand{\T}{\mathbb T}
\newcommand{\R}{\mathbb R}
\newcommand{\C}{\mathbb C}
\setlength{\parindent}{0pt}
\setlength{\parskip}{5pt}
\allowdisplaybreaks[1]
\begin{document}
\begin{center}
{\large\bfseries V1 frequency dispersion and excitation on a square torus}\\[4pt]
Explicit spectral derivations and interpretation of the simulations
\end{center}
This report extends the neural-field analysis in \texttt{final\_report.pdf}.
All new simulations use the three- and four-population calibrations specified
in Report 1. We distinguish exact statements, a local interval certificate,
ordinary numerical calculations, and a formal weak-forcing expansion.

\section{Field equation and the dispersion matrix}
Let $u:\T_L^d\times\R_+\to\R^N$, with $N=3,4$ and $d=1,2$, satisfy
\begin{equation}
 D\partial_tu=-u+W*f(u)+I(x,t),\qquad
 D=\diag(\tau_1,\ldots,\tau_N),\quad f(0)=0.
 \label{eq:field}
\end{equation}
The convolution applies the scalar nonlinearity to each \emph{source} population
before integration. Rows index targets and columns index sources. Our kernels are
\begin{equation}
 W_{ij}(x;p)=J_{ij}(p)K_{\Sigma_{ij}}(x),\quad
 K_\sigma(x)=\sum_{r\in\mathbb Z^d}
 \frac{\exp(-|x+Lr|^2/(2\sigma))}{(2\pi\sigma)^{d/2}},\quad
 J(p)=(0.5+p)G.
 \label{eq:kernel}
\end{equation}
Each $K_\sigma$ has integral one over the torus. With
$\widehat u_k=L^{-d}\int_{[0,L]^d}u(x)\ee^{-\ii q_k\cdot x}dx$,
$q_k=2\pi k/L$, convolution has multiplier
\begin{equation}
 \widehat W(q;p)=J(p)\odot\exp(-\Sigma|q|^2/2).
 \label{eq:what}
\end{equation}
Here $\widehat W$ denotes the integral multiplier, not the normalized Fourier
coefficient of $W$, and $\odot$ and the exponential act entrywise.
Writing $S=\diag(s_j)$, $s_j=f'_j(0)$, and substituting
$u=v\ee^{\lambda t+\ii q_k\cdot x}$ into the linearization at $I=0,u=0$ gives
\begin{equation}
 \lambda Dv=-v+\widehat W(q_k;p)Sv,\qquad
 \boxed{M(q,p)=D^{-1}[-I_N+\widehat W(q;p)S]},\quad q=|q_k|.
 \label{eq:M}
\end{equation}
Thus $P(\lambda,q,p)=\det(\lambda I_N-M(q,p))=0$ is the dispersion relation.
For the sigmoid in Report 1, $s_j=g_j\ell(-g_j\theta_j)[1-\ell(-g_j\theta_j)]$.
Time is in ms; space is in $0.1$ mm units. Therefore
\begin{equation}
 q_{\rm phys}=10q\ \mathrm{rad/mm},\qquad
 f_{\rm Hz}=\frac{1000}{2\pi}\omega,\qquad \omega=\Ima\lambda.
\end{equation}
At the calibrated onset, $p_0=0.25$, $k_0=3$,
$q_0=6\pi/47.6023=0.3959799405$ and $\omega_0=0.08$ rad/ms.

\newpage
\section{What is the imaginary part of a dominant eigenvalue?}
Define $a(q)=\max_{\lambda\in\spec M(q)}\Rea\lambda$. When a unique real root,
or one conjugate pair, has this real part, define $\omega_d(q)=0$ or the
positive imaginary part of that pair, respectively. At a tie between unrelated
roots, there is no unique physical frequency. For a single plotted curve we use
\begin{equation}
 \omega_d(q)=\max\{\Ima\lambda:\lambda\in\spec M(q),\ \Rea\lambda=a(q)\}.
 \label{eq:dominance}
\end{equation}
This convention can be discontinuous. It is different from continuing the same
complex eigenvalue after it ceases to dominate.

\begin{proposition}[Exact frequency derivative]
On an interval where $\lambda(q)=\alpha(q)+\ii\omega(q)$ is a simple eigenvalue,
let $Mv=\lambda v$ and $w^*M=\lambda w^*$, normalized by $w^*v=1$.
Then
\begin{equation}
 \lambda'=w^*M'v=-\frac{P_q}{P_\lambda},\qquad
 M'=-qD^{-1}[(\Sigma\odot J\odot\ee^{-\Sigma q^2/2})S].
 \label{eq:derivative}
\end{equation}
If $P_\lambda=A+\ii B$ and $P_q=C+\ii E$, evaluated at the root, then
\begin{equation}
 \boxed{\omega'(q)=\frac{CB-EA}{A^2+B^2}}.
 \label{eq:sign}
\end{equation}
Consequently a simple dominant pair increases strictly on any interval where
$CB-EA>0$ throughout and dominance does not change.
\end{proposition}
\begin{proof}
Differentiate $Mv=\lambda v$ and multiply by $w^*$. The two $v'$ terms cancel.
Alternatively differentiate $P(\lambda(q),q)=0$; simplicity gives
$P_\lambda\ne0$. Multiplying $-P_q/P_\lambda$ by
$\overline{P_\lambda}/\overline{P_\lambda}$ gives \eqref{eq:sign}.
Positive derivative implies strict increase by integration.
\end{proof}
No sign follows merely from Gaussian attenuation. If
$R_{ij}=J_{ij}s_j\ee^{-\Sigma_{ij}q^2/2}/\tau_i$, then
\begin{equation}
 \omega'=-q\sum_{ij}\Sigma_{ij}R_{ij}\Ima(\overline w_i v_j).
 \label{eq:edges}
\end{equation}
Both the signed edges and the complex eigenvectors matter. Even if all absolute
couplings decrease, the imaginary part of their collective eigenvalue can grow.

For comparison, if all $\tau_i=\tau$ and all retained widths equal $\sigma$,
then $M=(-I+hJS)/\tau$, $h=\ee^{-\sigma q^2/2}$.
For every eigenvalue $\nu$ of $JS$,
\begin{equation}
 \lambda_\nu(q)=(-1+h\nu)/\tau,\qquad
 \omega_\nu'=-\sigma q\omega_\nu<0\quad(q>0,\ \omega_\nu>0).
\end{equation}
Thus connectivity topology alone cannot prove increasing frequency. Unequal
widths and time scales require the actual sign calculation in \eqref{eq:sign}.

\newpage
\section{Explicit characteristic polynomials for the V1 motif}
At each $(q,p)$ write the two matrices as
\begin{equation}
 M_3=\begin{pmatrix}a&b&c\\d&e&f\\g&0&h\end{pmatrix},\qquad
 M_4=\begin{pmatrix}a&b&c&0\\d&e&f&0\\g&0&h&j\\r&0&s&t\end{pmatrix}.
 \label{eq:sparse}
\end{equation}
Every entry is explicit from \eqref{eq:M}; in particular
\begin{equation}
 (M_{ij})'=-q\Sigma_{ij}(M_{ij}+\delta_{ij}/\tau_i).
 \label{eq:entries}
\end{equation}
Define $H=(\lambda-a)(\lambda-e)-bd$ and $R=c(\lambda-e)+bf$.
Expanding the $3\times3$ determinant yields
\begin{equation}
 P_3=(\lambda-h)H-gR=\lambda^3+A\lambda^2+B\lambda+C,
 \label{eq:cubic}
\end{equation}
where
\begin{align}
 A&=-a-e-h,\notag\\
 B&=ae-bd+h(a+e)-gc,\notag\\
 C&=-h(ae-bd)+gce-gbf. \label{eq:abc}
\end{align}
For four populations, eliminate the last component when $\lambda\ne t$:
the third row acquires $jr/(\lambda-t)$ and $js/(\lambda-t)$ in its first
and third entries. Multiplying the reduced determinant by $\lambda-t$ gives
\begin{equation}
 \boxed{P_4=(\lambda-t)P_3-j(sH+rR)}.
 \label{eq:quartic}
\end{equation}
Both sides are polynomials, so equality also holds at $\lambda=t$.
The $P_3$ in this identity is the four-population matrix's upper block; it
is not the independently recalibrated three-population model.
Writing $P_4=\lambda^4+c_1\lambda^3+c_2\lambda^2+c_3\lambda+c_4$ gives
\begin{align}
 c_1&=A-t,&c_2&=B-tA-js,\notag\\
 c_3&=C-tB+js(a+e)-jrc,\notag\\
 c_4&=-tC-js(ae-bd)+jrce-jrbf. \label{eq:quarticcoeff}
\end{align}
Equations \eqref{eq:entries}--\eqref{eq:quarticcoeff}, differentiated by the
product rule, provide every coefficient and its $q$ derivative directly from
$G,\Sigma,D,S$. No fitted polynomial or eigenvalue finite difference is needed.

For three populations with one real root $\rho$ and a complex pair,
factorization gives
\begin{align}
 \alpha&=-\tfrac12(A+\rho),&
 \omega^2&=B-2\alpha\rho-\alpha^2,\notag\\
 \rho'&=-\frac{A'\rho^2+B'\rho+C'}{3\rho^2+2A\rho+B},&
 \alpha'&=-\tfrac12(A'+\rho'),\notag\\
 \omega'&=\frac{B'-2(\rho+\alpha)\alpha'-2\alpha\rho'}{2\omega}.
 \label{eq:cubicfrequency}
\end{align}
The root $\rho$ is the unique real zero of \eqref{eq:cubic} in this case;
$\omega^2>0$ selects a pair and $\alpha>\rho$ checks its dominance.
At a Hopf point $\alpha=0$, $\omega^2=B$, $C=AB$, with $A>0$ for the remaining
root $-A$ to be stable.

\newpage
\section{Four-population frequency and the limit of monotonicity}
For $\lambda=\alpha+\ii\omega$, set $y=\omega^2>0$. The real part of $P_4=0$
and its imaginary part divided by $\omega$ are, explicitly,
\begin{align}
 F(\alpha,y,q)&=\alpha^4-6\alpha^2y+y^2+c_1(\alpha^3-3\alpha y)
                  +c_2(\alpha^2-y)+c_3\alpha+c_4=0,\notag\\
 H(\alpha,y,q)&=4\alpha^3+3c_1\alpha^2+2c_2\alpha+c_3-(4\alpha+c_1)y=0.
 \label{eq:realimag}
\end{align}
These equations determine the nonreal branches; dominance must additionally
be checked against the other roots. The required partial derivatives are
\begin{align}
 F_\alpha&=4\alpha^3-12\alpha y+3c_1\alpha^2-3c_1y+2c_2\alpha+c_3,\notag\\
 F_y&=-6\alpha^2+2y-3c_1\alpha-c_2,\notag\\
 H_\alpha&=12\alpha^2+6c_1\alpha+2c_2-4y,&H_y&=-(4\alpha+c_1),\notag\\
 F_q&=c_1'(\alpha^3-3\alpha y)+c_2'(\alpha^2-y)+c_3'\alpha+c_4',\notag\\
 H_q&=c_1'(3\alpha^2-y)+2c_2'\alpha+c_3'.
\end{align}
Implicit differentiation and Cramer's rule give the fully real formula
\begin{equation}
 \boxed{\omega'=\frac{H_\alpha F_q-F_\alpha H_q}
 {2\omega(F_\alpha H_y-F_yH_\alpha)}}.
 \label{eq:quarticfrequency}
\end{equation}
At a simple root,
$F_\alpha H_y-F_yH_\alpha=|P_\lambda|^2/(2\omega^2)>0$:
the Jacobian of $(\Rea P,\Ima P)$ with respect to $(\alpha,\omega)$ has determinant
$|P_\lambda|^2$ by the Cauchy--Riemann equations, while $\Ima P=\omega H$ and
$y=\omega^2$. Thus the numerator in \eqref{eq:quarticfrequency} is an exact
sign test. At Hopf, $c_3=c_1y$, $c_4=c_2y-y^2$, and
\begin{equation}
 P_4=(\lambda^2+y)[\lambda^2+c_1\lambda+(c_2-y)].
\end{equation}
The noncritical roots are stable precisely when $c_1>0$ and $c_2-y>0$.

\begin{proposition}[No global nondecreasing positive frequency]
If every retained $\Sigma_{ij}>0$ and $\tau_i>0$, then every eigenvalue's
imaginary part tends to zero as $q\to\infty$. Hence if $\omega_d(q_*)>0$,
$\omega_d$ cannot be nondecreasing on $[q_*,\infty)$, nor on all subsequent
torus wavenumbers. If the $\tau_i$ are distinct, all roots are real for
sufficiently large $q$.
\end{proposition}
\begin{proof}
$M(q)=-D^{-1}+E(q)$, where $\|E(q)\|\le C\ee^{-\sigma_*q^2/2}\to0$ and
$\sigma_*$ is the smallest retained variance. Continuity of polynomial roots
forces the spectrum to the real set $\{-1/\tau_i\}$. A nondecreasing function
positive somewhere cannot tend to zero. For distinct $\tau_i$, take disjoint
small disks about these real limiting roots. For large $q$, each contains
exactly one root counted with multiplicity. Complex conjugation preserves each
disk and the spectrum of a real matrix, so its single root must be real.
\end{proof}
This theorem does not prohibit increasing frequency over a finite range.

\newpage
\section{What actually increases in the simulations?}
The continuous curves interpolate the allowed wavenumbers $q_k=2\pi k/L$.
The following values are ordinary double-precision root calculations;
Section~\ref{sec:certificate} proves increase on a smaller interval.
At the common onset, \eqref{eq:derivative} gives
\begin{equation}
 \omega'_3(q_0)=0.03316726,\qquad
 \omega'_4(q_0)=0.04570080.
 \label{eq:slopes}
\end{equation}
These derivatives use model-space $q$; divide by ten for physical rad/mm.
The constraint $\alpha'(q_0)=0$ locates the peak of \emph{growth}, not the
peak of frequency. There is no reason for $\omega'(q_0)$ to vanish.

\begin{center}\small
\begin{tabular}{lrrr}\toprule
At $p=0.25$ & Pair becomes dominant & Frequency maximum & Pair loses dominance\\
&\multicolumn{3}{c}{Wavenumber in rad/mm}\\\midrule
Calibrated 3-pop.&0.634537&11.174748&17.232247\\
Calibrated 4-pop.&1.937968&No interior maximum&10.048565\\
\bottomrule\end{tabular}
\end{center}
The three-population pair rises to $\omega=0.0967921$ and then decreases,
remaining dominant until $\omega=0.0833387$ at the last switch.
The four-population pair rises to a left limit $0.1009536$ at its switch;
the plotted dominant frequency then drops to zero because a \emph{different,
real eigenvalue} becomes rightmost. This is not a smooth frequency maximum,
and the complex pair does not disappear there.

\begin{center}\small
\begin{tabular}{rrrr}\toprule
Mode $k$&Calibrated 3-pop.&Calibrated 4-pop.&Original 4-pop.\\\midrule
0&0&0&0\\1&0.073439&0&0.071180\\2&0.076059&0.074342&0.075302\\
3&0.080000&0.080000&0.081049\\4&0.084551&0.086069&0.087252\\
5&0.088966&0.091700&0.093016\\6&0.092675&0.096341&0.097763\\
7&0.095304&0.099657&0.101153\\8&0.096640&0&0.103010\\
9&0.096593&0&0\\10&0.095167&0&0\\\bottomrule
\end{tabular}
\end{center}
Entries are $\omega_d$ in rad/ms. The original four-population matrix is
evaluated at its own computed mode-3 onset, $p=0.2500000483$.
Its modes 1--8 do increase; it switches to a real dominant root at
$11.158365$ rad/mm, before mode 9. Thus the original increasing plot is
consistent with the analysis, but does not prove global monotonicity.

Equation~\eqref{eq:edges} also explains the sign quantitatively. At $q_0$, the
positive edge contributions sum to $0.11992572$ (3-pop.) and $0.15464963$
(4-pop.); negative contributions sum to $-0.08675846$ and $-0.10894883$.
For example, attenuation of E$\to$E contributes $+0.08898428$ and $+0.11141269$,
whereas E$\to$PV contributes $-0.03193637$ and $-0.04360907$.
The sums give \eqref{eq:slopes}. These are modal sensitivities, not increases
in the physical coupling magnitudes.

\newpage
\begin{figure}[H]
\centering\includegraphics[width=\textwidth,height=.78\textheight,keepaspectratio]{dispersion_branches.pdf}
\caption{The left panels compare the continued complex pair's real part with
the largest real eigenvalue. The right panels distinguish the continued pair's
frequency (gray) from the dominant frequency (blue); points mark allowed torus
modes. Vertical lines locate dominance switches. The onset is not a frequency
maximum. In four populations, the abrupt drop is caused by a real root taking
over while the complex pair still oscillates.}
\label{fig:branches}
\end{figure}
The behavior persists in the sampled parameter cases $p=0.23,0.2495,0.25,0.27$.
For three populations, the smooth frequency peak moves from $11.1012$ to
$11.2445$ rad/mm; for four, the upper dominance switch moves from $9.4203$ to
$10.6089$ rad/mm. These sampled continuations are numerical observations,
not interval proofs over the entire $(p,q)$ rectangle.

\newpage
\section{A certified local increase for both V1 models}\label{sec:certificate}
The derivative at a single point implies increase on some neighborhood by
continuity, provided the root is simple and remains dominant. To give an
\emph{explicit} neighborhood without relying on a plotting grid, the supplied
interval calculation uses $p=1/4$ and
\begin{equation}
 Q=[0.39597,0.39599],\qquad Q_{\rm phys}=[3.9597,3.9599]\ \mathrm{rad/mm}.
 \label{eq:Q}
\end{equation}
This interval contains $q_0$. Model constants are interpreted as the exact
decimal strings saved in \texttt{calibration.json}; the result is conditional
on those coefficients, not a certificate about uncertain biological parameters.

Here is the validation argument. Set
$\mathcal F(x,q)=(\Rea P(\alpha+\ii\omega,q),\Ima P(\alpha+\ii\omega,q))$
with $x=(\alpha,\omega)$. Choose a numerical center $x_0$, a rectangle
$X=x_0+[-10^{-4},10^{-4}]^2$, and a constant nonsingular approximation
$C$ to $\mathcal F_x(x_0,q_0)^{-1}$. Interval arithmetic encloses
\begin{equation}
 E=I-C\mathcal F_x(X,Q),\qquad
 \mathcal K=x_0-C\mathcal F(x_0,Q)+E(X-x_0).
 \label{eq:krawczyk}
\end{equation}
For each fixed $q\in Q$, the mean-value theorem gives
$x-C\mathcal F(x,q)\in\mathcal K$ for every $x\in X$.
If $\mathcal K\subset\operatorname{int}X$ and $\|E\|_\infty<1$, this map
is a contraction of $X$ into itself. Banach's theorem yields exactly one root
in $X$ for every $q\in Q$. The same scalar construction encloses every real
root. The conjugate rectangle and the disjoint real intervals account for all
$N$ polynomial roots. Since their real parts are below $\inf X_\alpha$, the
complex pair is dominant throughout $Q$. Invertibility of the Jacobian gives
simplicity and a differentiable root branch.

The implementation uses outward-directed 60-digit decimal arithmetic, expands
each correctly rounded exponential to its two neighboring decimal values, and
evaluates \eqref{eq:sign} on the entire root rectangle and $Q$.
The following are conservative outward roundings of its bounds:
\begin{center}
\begin{tabular}{lrr}\toprule
Certified quantity&3 populations&4 populations\\\midrule
Contraction bound $\|E\|_\infty$&$<0.00910$&$<0.02001$\\
Lower bound for $\omega'$&0.03259&0.04407\\
Upper bound for $\omega'$&0.03375&0.04739\\
Other real roots&$[-.013984,-.013783]$&$[-.039901,-.039700]$\\
&&$[-.001982,-.001781]$\\\bottomrule
\end{tabular}
\end{center}
The positive lower bounds prove that \emph{both calibrated models have a
strictly increasing dominant frequency on \eqref{eq:Q}}.
They do not certify the much wider increasing ranges suggested by Figure~\ref{fig:branches}.
The certificate is reproducible with
\texttt{python -m simulations.v1\_interval\_certificate}; its full rectangles,
Newton images and bounds are saved in \texttt{local\_monotonicity\_certificate.json}.

\newpage
\section{The square torus: modes and an exact non-excitation theorem}
For $d=2$, write
\begin{equation}
 u(x,y,t)=\sum_{m,n\in\mathbb Z}\widehat u_{mn}(t)
 \ee^{2\pi\ii(mx+ny)/L},\qquad
 q_{mn}=\frac{2\pi}{L}(m,n).
\end{equation}
The linearized forced equation is
\begin{equation}
 \partial_t\widehat u_{mn}=M(|q_{mn}|,p)\widehat u_{mn}+D^{-1}\widehat I_{mn}.
 \label{eq:fourier2}
\end{equation}
It does not mix different $(m,n)$, although it mixes populations. A square
lattice has discrete rotational/reflection symmetry; radial kernels do not
make its allowed directions continuously rotationally invariant.
At the selected radius, $m^2+n^2=9$ has exactly the four integer solutions
$(\pm3,0),(0,\pm3)$. Thus the $+\ii\omega_0$ eigenspace has four complex
amplitudes, and the real center space has dimension eight. The computed
shell check finds all other shells stable at onset. Degeneracy of the x and y
modes does not itself couple them.

\begin{proposition}[An exact axis remains an axis]\label{prop:axis}
Assume the kernels are translation invariant, $f$ is globally Lipschitz, and
$I(x,y,t)=I(x,t)$, $u(x,y,0)=u_0(x)$. Then the unique solution of
\eqref{eq:field} is independent of $y$ for every $t\ge0$. In particular,
$\widehat u_{mn}(t)=0$ for every $n\ne0$, even with the full nonlinearity.
\end{proposition}
\begin{proof}
On $C(\T_L^2;\R^N)$ the right-hand side is globally Lipschitz in $u$:
convolution is bounded by the kernels' $L^1$ norms and the componentwise
sigmoids have Lipschitz constants $g_j/4$. Standard Picard iteration yields
uniqueness and global existence for continuous forcing bounded on finite time
intervals. For any shift $h$, $v(x,y,t)=u(x,y+h,t)$ solves the same equation,
with the same forcing and initial condition. Uniqueness gives $v=u$ for every
$h$, proving independence of $y$ and hence the Fourier assertion.
\end{proof}
The statement includes a traveling axial wave $I\propto\cos(qx-\Omega t)$,
a static stripe $I\propto\cos(qx)$, and their temporal modulations.
Nonlinearity can generate higher x harmonics, but cannot produce transverse
modes from exactly zero. Indeed multiplication adds wavevectors,
\begin{equation}
 \widehat{u_j^2}_{mn}=\sum_{a,b}\widehat u_{j,ab}\widehat u_{j,m-a,n-b};
 \qquad (a,0)+(b,0)=(a+b,0).
\end{equation}
The translation proof also covers non-polynomial nonlinearities without
requiring convergence of a global Taylor series. More generally Fourier
support in an additive lattice subgroup is preserved when the initial data
and forcing share that support (equivalently their translation symmetries).

\newpage
\section{How transverse modes can grow when they are seeded}
Proposition~\ref{prop:axis} is an invariance theorem, not a stability theorem.
Let $U(x,t)$ be an axial solution and put $u=U+v$. To first order,
\begin{equation}
 Dv_t=-v+W*[f'(U)v].
 \label{eq:transverse}
\end{equation}
If $a_{j,\ell}(t)$ is the x-Fourier coefficient of $f'_j(U_j(x,t))$, then
\begin{equation}
 \dot v_{i,mn}=-\frac{v_{i,mn}}{\tau_i}
 +\frac1{\tau_i}\sum_j\widehat W_{ij}(|q_{mn}|)
       \sum_{\ell\in\mathbb Z}a_{j,m-\ell}(t)v_{j,\ell n}.
 \label{eq:transversefourier}
\end{equation}
Different x wavenumbers can now couple within each fixed transverse index $n$.
There is still no inhomogeneous source: $v_{\cdot,\cdot n}=0$ remains zero.
For a periodic axial solution, stability of a finite Fourier truncation is
determined by its monodromy matrix $\mathcal U_n(T,0)$:
\begin{equation}
 \gamma_n=T^{-1}\log\rho(\mathcal U_n(T,0)).
 \label{eq:floquet}
\end{equation}
A multiplier of modulus greater than one amplifies some seeded transverse
perturbation; multipliers strictly inside the unit circle imply linear decay
for that truncation. Equation~\eqref{eq:fourier2} alone is insufficient around
a spatially patterned or time-dependent $U$. For the full function-space
problem one needs the corresponding evolution-operator analysis and control
of the Fourier truncation; a finite-time trace is not such a proof.

\subsection{Weak-forcing selection rules, derived by expansion}
These rules explain which amplitudes can couple at leading order. They are a
formal multiple-scale approximation, separate from the exact theorem above.
Set $p=p_0+\varepsilon^2\mu$, $T=\varepsilon^2t$,
$I=\varepsilon^2F_2$, and $u=\varepsilon u_1+\varepsilon^2u_2+\cdots$.
Let $M(q_0)\Phi=\ii\omega_0\Phi$, $\Psi^*M(q_0)=\ii\omega_0\Psi^*$,
$\Psi^*\Phi=1$. With $q_x=(q_0,0)$, $q_y=(0,q_0)$,
\begin{equation}
 u_1=\sum_{\nu\in\{\pm q_x,\pm q_y\}}A_\nu(T)\Phi
       \ee^{\ii(\omega_0t+\nu\cdot x)}+\text{c.c.}
 \label{eq:criticalexpansion}
\end{equation}
Define the quadratic multiplier and nonresonant resolvent by
\begin{equation}
 Q_k(v,w)=\tfrac12D^{-1}\widehat W(k)
                  [f''(0)\odot v\odot w],\qquad
 R(\Omega,k)=(\ii\Omega I_N-M(|k|))^{-1}.
\end{equation}
A forcing harmonic $F_{\Omega,s}\ee^{\ii(\Omega t+s\cdot x)}$ contributes
$U_{\Omega,s}=R(\Omega,s)D^{-1}F_{\Omega,s}$ to $u_2$ when the inverse exists.
At order $\varepsilon^3$, the term $2Q(u_1,u_2)$ is projected by $\Psi^*$
onto each critical harmonic. Matching both frequency and wavevector gives the
rules on the next page; the factor two comes from the two quadratic orderings.

\newpage
\subsection{Temporal, spatial and traveling pumps}
For a uniform temporal pump at $2\omega_0$, the product of
$U_{2\omega_0,0}\ee^{2\ii\omega_0t}$ and
$\overline{A_\mu\Phi}\ee^{-\ii(\omega_0t+\mu\cdot x)}$
resonates with mode $\nu$ precisely when $\nu=-\mu$. Hence
\begin{equation}
 (\partial_T A_\nu)_{\rm temporal}
 =\eta_t\overline{A_{-\nu}},\qquad
 \eta_t=2\Psi^*Q_\nu(\overline\Phi,U_{2\omega_0,0}).
 \label{eq:temporalselection}
\end{equation}
The same coefficient applies to both axes by symmetry. This is parametric
amplification, not an additive source for a previously absent spatial mode.

For a static spatial pump at $s=2q_x$, the product of its second-order response
and $A_\mu\Phi\ee^{\ii(\omega_0t+\mu\cdot x)}$ resonates when $\nu-\mu=s$.
Among the critical modes, only $(\nu,\mu)=(q_x,-q_x)$ satisfies this:
\begin{equation}
 (\partial_T A_{q_x})_{\rm spatial}=\eta_x A_{-q_x},\qquad
 \eta_x=2\Psi^*Q_{q_x}(\Phi,U_{0,2q_x}).
 \label{eq:spatialselection}
\end{equation}
The $-2q_x$ harmonic supplies the reverse coupling. For a real cosine the
two coefficients coincide, because the static response is real and kernels
are even. There is no such leading linear coupling into $A_{\pm q_y}$.
An oblique pump $s=q_y-q_x$ could instead couple the two axes; it is not an
axial stimulus.

For a traveling pump $(\Omega,s)=(2\omega_0,2q_x)$, resonance with a conjugate
critical mode requires $\nu+\mu=2q_x$. Since $|\nu|=|\mu|=q_0$, equality in
the triangle inequality forces $\nu=\mu=q_x$. Therefore
\begin{equation}
 (\partial_T A_{q_x})_{\rm traveling}
 =\eta_w\overline{A_{q_x}},\qquad
 \eta_w=2\Psi^*Q_{q_x}(\overline\Phi,U_{2\omega_0,2q_x}).
\end{equation}
The opposite traveling mode and the transverse pair are not directly resonant
with this single traveling harmonic at this order.
Reversing propagation exchanges $q_x$ and $-q_x$; the notebook uses
$\cos(2q_x\cdot x-2\omega_0t)$.

Nonlinear amplitude interactions can nevertheless change transverse stability.
For example, square symmetry permits terms
\begin{equation}
 \partial_T B_\pm=\gamma(T)B_\pm+\kappa(T)\overline{B_\mp}+O(|B|^2),
 \quad \gamma=a+d(|A_+|^2+|A_-|^2),\quad
 \kappa=eA_+A_-+\eta_t,
\end{equation}
where $A_\pm=A_{\pm q_x}$ and $B_\pm=A_{\pm q_y}$; $a,d,e$ are the
projected normal-form coefficients. These terms vanish at $B=0$, as the exact
theorem requires. If $\gamma,\kappa$ are constant in the chosen rotating frame,
the matrix acting on $(B_+,\overline B_-)$ and its exponents are
\begin{equation}
 \begin{pmatrix}\gamma&\kappa\\\overline\kappa&\overline\gamma\end{pmatrix},
 \qquad \Lambda_\pm=\Rea\gamma\pm\sqrt{|\kappa|^2-(\Ima\gamma)^2}.
\end{equation}
Positive real part permits seeded growth; zero initial $B$ still stays zero.
For varying coefficients one must instead solve the variational equation.

\newpage
\section{Direct simulations, frequency interpretation and reproducibility}
The notebook integrates \eqref{eq:field} using Fourier convolution and RK4,
with the coefficients in Report 1. The 1D runs use 64 points, a 1 ms step,
80 s duration and saved output every 10 ms. At $p=0.2495$, homogeneous temporal
forcing at $2\omega_0$ or static spatial forcing at mode 6 each produces an
oscillatory mode-3 response above the decaying unforced control. The table
averages $|\Psi^*\widehat u_{k_0}|$ over the final 16 s:
\begin{center}\small
\begin{tabular}{lrr}\toprule
Late projected mode-3 magnitude&3 populations&4 populations\\\midrule
Unforced&$2.58\times10^{-5}$&$1.47\times10^{-5}$\\
Temporal only&$1.84\times10^{-3}$&$1.37\times10^{-3}$\\
Static spatial only&$2.23\times10^{-3}$&$1.74\times10^{-3}$\\\bottomrule
\end{tabular}
\end{center}
The resolved late response is about $12.75$ Hz, near $1000\omega_0/(2\pi)=12.7324$
Hz. Temporal forcing is at $25.4648$ Hz, so this is
the predicted near-subharmonic response. Static forcing has zero temporal
frequency; its oscillatory response is generated by the field dynamics.
The late FFT resolution is $0.0625$ Hz (16 s window).
The dispersion curve $\omega_d(q,p)$ describes the unforced linearization
around zero. It is not the forced oscillation frequency for every nonlinear
state. Damped eigenfrequencies need not appear as persistent peaks in a
forced or noisy trace, and branch switching does not require a visible
nonlinear oscillation to stop abruptly.

The 2D runs use a $32\times32$ square torus, 1 ms steps and 60 s duration.
The reported transverse norm is the RMS over populations and grid points of
$u-\langle u\rangle_y$. Exactly axial initial data remain axial. With a
nonzero transverse seed, the outcomes are:
\begin{center}\small
\begin{tabular}{lrr}\toprule
Transverse RMS&3 populations&4 populations\\\midrule
Initial seeded value&$1.225\times10^{-4}$&$1.061\times10^{-4}$\\
Final, axial traveling drive&$1.327\times10^{-6}$&$3.888\times10^{-8}$\\
Final, axial static drive&$9.949\times10^{-7}$&$7.064\times10^{-7}$\\
Final, uniform temporal drive&$2.298\times10^{-3}$&$1.009\times10^{-3}$\\
Final, unforced&$1.522\times10^{-5}$&$6.859\times10^{-6}$\\\bottomrule
\end{tabular}
\end{center}
These examples demonstrate decay under the tested axial drives and growth
under uniform temporal pumping. They neither exclude transverse instability
under other axial drives nor estimate a converged Floquet exponent. The exact
answer is: an axis wave cannot create transverse modes from perfectly axial
data; it may change the growth of transverse modes already present.

Validation checks preserve the connectivity masks, compare sparse determinants
and eigenvalue derivatives, verify critical square-lattice shells and the
Gaussian high-wavenumber stability tail, and compare time steps.
A $32$ versus $48$ grid comparison at 2 s for the 2D spatial case has relative
errors $1.41\times10^{-5}$ and $3.42\times10^{-5}$; this short check is not
an error bound for the entire 60 s trajectory.

\begingroup\small\raggedright
\textbf{Reproduce.} The executed notebook is
\path{simulations/v1_three_four_population.ipynb}.
The full coefficient and simulation records are in
\path{simulations/output/v1_extension/}.
\texttt{v1\_dispersion\_detail.py} generates the branch figure and switch data;
\texttt{v1\_interval\_certificate.py} validates the local increase.
The LaTeX sources and PDF builder are in \texttt{notes/}. Report 1 documents
the biological references; all spectral and symmetry results here are derived
from the explicitly stated field equation.
\par\endgroup
\end{document}
